\subsubsection{Microbial turnover across Arctic and North Atlantic environmental gradients}

Marine microorganisms sustain major ocean biogeochemical cycles, yet their
community composition varies across water masses, depth layers and
physicochemical gradients~\cite{sunagawa2015structure,salazar2019gene}. We
reanalysed ribosomal-RNA miTAG community profiles and matched environmental
metadata from the Tara Oceans expedition, which sampled planktonic communities
across the global ocean between 2009 and
2013~\cite{pesant2015open,salazar2019gene}. The source dataset contained 180
prokaryote-enriched metagenomic samples. We selected 38 Arctic Ocean samples
to represent an environmentally extreme and comparatively understudied polar
ecosystem relative to lower-latitude oceans. We selected 24 North Atlantic
Ocean samples because their spatial distribution broadly traces major
circulation pathways, which can delimit plankton genomic
provinces~\cite{richter2022currents}. Together, the selected samples span
surface, deep-chlorophyll-maximum, mixed and mesopelagic waters
(Fig.~\ref{fig:ocean-hric-regions}). The mapped currents provide geographical
context only; current transport was not included as a covariate in the present
analysis.

We treated each water sample as a local community and each ocean region as its
regional pool. HRIC alpha and gamma diversity were therefore calculated at the
sample and region levels, respectively, and sample-level turnover was defined
as $\gamma-\alpha$. Mean alpha diversity was 0.270 in the Arctic and 0.368 in
the North Atlantic, whereas regional gamma diversity was 0.379 and 0.503,
respectively. The corresponding additive beta diversities, equal to mean
turnover within each region, were 0.109 and 0.135. These values describe the
two selected regional pools; they are not a formal test of an Arctic--North
Atlantic difference. Ordering samples from the smallest to the largest
turnover linked each local diversity value to the taxonomic composition shown
directly below it.

Two environmental trends were consistent across regions. Turnover declined
with depth in both the Arctic (Spearman's $\rho=-0.68$, $P=5.1\times10^{-6}$,
$n=36$) and North Atlantic ($\rho=-0.60$, $P=0.0017$, $n=24$), and increased
with chlorophyll $a$ (Arctic, $\rho=0.64$, $P=1.7\times10^{-5}$, $n=38$;
North Atlantic, $\rho=0.68$, $P=2.5\times10^{-4}$, $n=24$). The Arctic
gradient also included higher turnover at greater oxygen concentration
($\rho=0.61$, $P=4.6\times10^{-5}$, $n=38$) and lower turnover at greater
phosphate ($\rho=-0.59$, $P=0.0010$, $n=28$) and nitrate concentrations
($\rho=-0.66$, $P=1.2\times10^{-4}$, $n=28$). Arctic turnover was not
associated with temperature ($\rho=-0.08$, $P=0.65$, $n=38$). In the North
Atlantic, turnover increased with temperature ($\rho=0.44$, $P=0.033$,
$n=24$) and declined with nitrate ($\rho=-0.48$, $P=0.023$, $n=22$).
Associations with oxygen ($P=0.051$) and phosphate ($P=0.33$) were not
supported at the 0.05 level. These are exploratory, univariable correlations
with unadjusted two-sided $P$ values. Depth, temperature, oxygen, chlorophyll
and nutrient concentrations covary in the water column, so their coefficients
should be interpreted as environmental gradients rather than independent
effects.

The stacked profiles showed that samples with similar turnover could contain
different dominant taxa, and that the taxonomic signature of the environmental
gradient differed between regions. We used nitrate for the taxon-resolved
panel because, among the six displayed covariates, it yielded the largest
number of significant univariable associations with HRIC-transformed taxon
coordinates. This choice was data-driven, so the taxon-level results are
exploratory. Among the 19 most abundant taxa in each region, 14 Arctic and 13
North Atlantic taxa were associated with nitrate after Benjamini--Hochberg
adjustment within region ($q<0.05$). In the Arctic, the HRIC coordinates of
SAR86, SAR11 and Marinimicrobia increased with nitrate, whereas those of
\textit{Formosa}, Cryomorphaceae, NS5, \textit{Polaribacter}, SAR92,
\textit{Marinoscillum} and Flavobacteriaceae decreased. In the North Atlantic,
SAR324, SAR202 and Marinimicrobia increased with nitrate, whereas NS4, OM1,
Gammaproteobacteria, \textit{Prochlorococcus}, NS5, SAR86 and
\textit{Marinoscillum} decreased. Marinimicrobia showed the same positive
association in both regions, and NS5, chloroplast-assigned sequences and
\textit{Marinoscillum} showed negative associations in both. By contrast,
SAR86 changed sign between regions. These coefficients describe changes in
HRIC taxon coordinates, not changes in relative abundance alone, and do not by
themselves identify causal drivers of turnover.

Together, the regional analyses identify a shared depth--productivity axis
along which microbial turnover changed in both oceans, superimposed on
region-specific temperature, oxygen and nutrient patterns. The taxon-resolved
results further show that a similar nitrate--turnover trend need not arise from
the same community reorganization in each basin. Because the samples are
observational, unevenly distributed among water-column layers and exposed to
correlated environmental gradients, these results support environmental
filtering as a hypothesis rather than a causal conclusion. They nevertheless
demonstrate how local diversity, regional diversity, turnover and individual
taxon coordinates can be followed in one coherent analysis.

\begin{figure*}[t]
    \centering
    \includegraphics[
        width=\textwidth,
        height=0.82\textheight,
        keepaspectratio
    ]{figs/figure_integrated_ocean_regions.pdf}

    \caption{\textbf{Environmental gradients accompany microbial turnover
    across Arctic and North Atlantic waters.}
    \textbf{a--d,} Arctic Ocean communities (\(n=38\)).
    \textbf{e--h,} North Atlantic Ocean communities (\(n=24\)).
    \textbf{a,e,} Sampling locations, shown as black crosses. Schematic warm
    and cold currents in \textbf{e} provide geographical context and were not
    included in the statistical models.\ \textbf{b,f,} Sample-level alpha
    diversity calculated using \texttt{HRIC::SHalpha}. Samples are ordered
    from the smallest to the largest turnover, defined as
    \(\gamma-\alpha\); dashed red lines show region-specific gamma diversity
    calculated using \texttt{HRIC::SHgamma}.\ \textbf{c,g,} Relative abundance
    of the 19 taxa with the largest mean relative abundance in each region,
    with all remaining taxa grouped as Others. Bars follow the sample order in
    \textbf{b} and \textbf{f}. Asterisks identify taxa whose HRIC-transformed
    coordinate was associated with nitrate in univariable linear models at
    Benjamini--Hochberg-adjusted \(q<0.05\), with adjustment across the 19
    tested taxa within each region.\ \textbf{d,h,} Turnover in relation to
    depth, temperature, oxygen, chlorophyll \(a\), phosphate and nitrate. Red
    lines and shaded areas are ordinary least-squares visual summaries with
    95\% confidence intervals. Annotations report Spearman's \(\rho\),
    unadjusted two-sided \(P\) values and complete-case sample sizes.}

    \label{fig:ocean-hric-regions}
\end{figure*}
